% ============================================================================
%  SECTION: Bring-Up Failures, Modeling-Driven Diagnosis, and Fixes
%  Status: MOCK PROSE / PLANNING DRAFT. Numbers are traceable to
%  docs/boost-bringup-debug.md and tools/inductance/ results, but several are
%  flagged UNCONFIRMED in the source log and must not be hardened into claims
%  before the owed bench measurements land.
% ============================================================================

\section{Bring-Up Failures, Modeling-Driven Diagnosis, and Fixes}
\label{sec:bringup}

% NOTE: framing paragraph. The reviewer's likely reaction to a "we broke five
% converters" section is "why is this in a control paper?" -- answer that in the
% first sentence, not the last. The thesis-level claim is that the droop
% architecture was never implicated, and that establishing that required a
% modeling campaign of the same character as the control synthesis itself.

The power-sharing controller developed in Section~\ref{sec:synthesis} was designed
against a small-signal plant that presumes a functioning pair of boost converters
delivering current into a common DC bus. Reaching that operating point on hardware
proved substantially harder than reaching it in simulation: five TPS61288 boost
stages were destroyed over a six-week bring-up campaign, none of them while the
droop loop was closed. This section reports that campaign because its outcome is
methodologically relevant rather than merely anecdotal. Two of the three root
causes were identified only after building quantitative models---a Gerber-derived
parasitic-inductance extraction and a soft-start charge-transfer timing model---and
both models resolved questions that no amount of additional bench iteration had
resolved in four prior attempts. We also report the hypotheses that were wrong,
and the evidence that superseded them, because the discarded hypotheses were each
individually plausible and are each individually the standard textbook answer.

% ------------------------------------------------------------------
\subsection{Chronology of the Converter Failures}
\label{sec:bringup:chronology}

% NOTE: keep this table close to the original log's Failure Datapoints table.
% The key rhetorical move is that the *variance* across deaths is what kills
% each hypothesis -- 120 mA and 5 A both fatal, battery and bench supply both
% fatal, original and reworked parts both fatal.

All five failures shared an identical post-mortem signature: the converter's input,
switch, and output nodes fused to ground ($V_\mathrm{BT}\!\rightarrow\!\mathrm{GND}$
or $V_\mathrm{FC}\!\rightarrow\!\mathrm{GND}$ short), consistent with destruction of
the synchronous-rectifier FET. Table~\ref{tab:deaths} summarizes the conditions.

\begin{table}[H]
\caption{Converter failure events during bring-up. Deaths 1--4 are battery-channel;
Death~5 is fuel-cell-channel.\label{tab:deaths}}
\begin{adjustwidth}{-\extralength}{0cm}
\centering
\begin{tabular}{clll}
\toprule
\textbf{\#} & \textbf{Source} & \textbf{Action} & \textbf{Outcome} \\
\midrule
1 & Two \SI{9}{\volt} batteries & Enable BT$\rightarrow$bus ideal-diode switch & BT boost destroyed (original, never-reworked part) \\
2 & Bench supply, \SI{120}{\milli\ampere} limit & Firmware bring-up: switches, then boosts & Supply collapsed into current limit; BT boost destroyed \\
3 & Bench supply, $\geq$\SI{5}{\ampere} & Manual bring-up command & \SIrange{5}{7}{\ampere} with collapsed rail; BT boost destroyed \\
4 & Bench supply, \SI{8.2}{\volt}/\SI{200}{\milli\ampere} & Known-good FC part transplanted to BT pad & Immediate current-limit peg; BT boost destroyed \\
5 & FC from stiff bench supply, VESC attached & Close motor-node switch at full bus & FC boost destroyed \\
\bottomrule
\end{tabular}
\end{adjustwidth}
\end{table}

% NOTE: this paragraph is the "honest scientific narrative" the section is
% supposed to deliver. Do not soften it in revision -- the superseded
% hypotheses are the most transferable content here.

Four hypotheses were advanced and subsequently rejected. \emph{(i) Capacitive
inrush.} The obvious explanation for a converter dying at the instant it is
connected to a bus is inrush into bulk capacitance. This was refuted
topologically: the \SI{470}{\micro\farad} bulk capacitor sits on the motor node
\emph{behind} a separate ideal-diode switch, not on the shared bus, and that
switch was open in every one of Deaths~1--4. The bus itself carries only the
\SIrange{30}{40}{\micro\farad} of ceramics distributed across the ideal-diode
controllers. \emph{(ii) Source collapse and restart cycling.} Because the logic
rail is derived from the battery input through a linear regulator, a soft source
sags under converter load, browns out the microcontroller, and provokes a
re-enable cycle---a genuine failure mode observed on the bench. It was rejected as
\emph{the} cause because Death~3 occurred from a stiff $\geq$\SI{5}{\ampere}
supply that never collapsed, and Death~2 occurred at \SI{120}{\milli\ampere},
where the deliverable input energy is orders of magnitude below the destruction
energy. \emph{(iii) Reverse conduction through the disabled converter.} An early
design note asserted that a disabled boost presents a body-diode path through
which regenerative energy back-feeds the synchronous rectifier. Reading the
RT1987 ideal-diode datasheet retired this: the part uses back-to-back integrated
FETs and provides full isolation when disabled. The claim survived in the design
documentation for several weeks after it was falsified, and propagated into
firmware sequencing comments---a documentation-hygiene failure worth reporting.
\emph{(iv) Compensation-network mismatch.} The battery channel's compensation
resistor differed from the fuel-cell channel's; matching them was the intervention
in Death~4, which still destroyed the part.

Death~4 is the decisive experiment. A converter that had been \emph{proven good}
regulating on the fuel-cell pad was transplanted to the battery pad, verified to
regulate standalone, and then destroyed on its first bus connection. Combined with
three non-destructive control experiments---driving the bus node directly from a
supply with the converter removed (clean, $\approx$\SI{0}{\ampere}); driving the
bus from the fuel-cell converter (clean); and OR-ing a live bus against a supply
through the battery-side ideal diode (clean)---the fault was localized to the
\emph{battery channel's physical copper}, with the schematic confirmed identical
part-for-part. That conclusion, however, was reached by elimination, and
elimination does not produce a mechanism. Producing one required the modeling
described next.

% ------------------------------------------------------------------
\subsection{Parasitic-Inductance Extraction from Fabrication Gerbers}
\label{sec:bringup:inductance}

% NOTE: this is the subsection most likely to interest an Energies reviewer,
% because it is a reusable method rather than a board-specific war story.
% Consider promoting it to its own top-level section if the paper runs long
% in the control chapters.

The surviving hypothesis was that the battery channel's output-capacitor
commutation loop---the ``hot loop'' carrying the discontinuous current pulse when
the synchronous rectifier commutates---enclosed enough area to ring the switch node
past the device's \SI{20}{\volt} absolute-maximum rating under bus-drive
$\mathrm{d}i/\mathrm{d}t$, with only \SIrange{0.5}{1.5}{\volt} of nominal headroom
available. Physical measurement on the assembled board found the battery channel's
output ceramics placed \SI{240}{mil} from the converter output pin against
\SI{40}{mil} on the fuel-cell channel. A placement ratio, however, is not an
inductance ratio: all power nets on this board are polygon pours, so conductor
\emph{length} is a poor proxy and an earlier trace-length-based estimate was later
withdrawn as an artifact of a parser that excluded the pours.

We therefore built a dedicated extraction tool
(\texttt{tools/inductance/}) that computes loop inductance directly from the
fabrication data, with no manual geometry re-entry:

\begin{enumerate}
\item \textbf{Gerber and Excellon parsing.} A self-contained RS-274X and drill
parser converts the manufacturing files into copper polygons and hole locations,
so the extracted geometry is the geometry that was actually fabricated.
\item \textbf{Net isolation by flood fill.} Gerbers carry no net names; the forward
net (a converter-output pour) and the return net (the ground pour on the opposite
layer) are selected by coordinate and flood-filled at a fine isolation pitch, then
resampled to a coarser meshing pitch. Plane-clearance slots are morphologically
bridged so the return plane remains a single connected body at every pitch, with a
severed-plane detector that flags meshing artifacts.
\item \textbf{PEEC solution.} The meshed forward and return conductors are emitted
as a uniform-mesh FastHenry deck at their true stackup separation, with an explicit
port and an explicit loop closure at the load, and solved for
$L(f) = \operatorname{Im}\{Z_\mathrm{port}(f)\}/2\pi f$ over DC to \SI{1}{\mega\hertz}.
\item \textbf{Independent bound.} A closed-form microstrip estimate is computed from
the same mesh as a sanity check on the numerical result.
\end{enumerate}

% NOTE: emphasize the sweep design -- it is not a cross product, and that is a
% defensible methodological choice worth one sentence in a journal paper.

Because a PEEC result is only as trustworthy as its discretization, the tool solves
two orthogonal sweeps rather than a cross product: a \emph{mesh-convergence} axis
(all mesh pitches at one designated frequency) and a \emph{frequency} axis (all
frequencies at the converged pitch), sharing the designated point. Solutions are
persisted per-point to an append-as-you-go store, making multi-hour runs
interruptible and idempotent.

The result is given in Table~\ref{tab:Lsweep}. At the converged mesh, the
battery-channel output loop extracts to \SI{3.48}{\nano\henry} at DC falling to
\SI{3.15}{\nano\henry} at \SI{1}{\mega\hertz} (skin effect), against
\SI{1.54}{\nano\henry} and \SI{1.46}{\nano\henry} for the fuel-cell channel---a
ratio of $2.2\times$ at \SI{1}{\mega\hertz}. The two channels' extracted loop
\emph{lengths} (\SI{9.69}{\milli\meter} versus \SI{4.74}{\milli\meter}) and
effective widths (\SI{5.91}{\milli\meter} versus \SI{12.47}{\milli\meter}) show the
asymmetry is both a longer and a narrower return path.

% TODO(verify): CLAUDE.md and the debug log quote "~2.7x"; the current
% freq_sweep store gives 2.16x @ 1 MHz / 2.26x @ DC. The 2.7 figure predates a
% change in the loop-closure definition (pre-2026-07-01 L numbers are stale).
% Decide ONE number for the paper, cite the store that produced it, and delete
% the other everywhere. Do not publish both.

\begin{table}[H]
\caption{Extracted output-loop inductance, converged mesh.\label{tab:Lsweep}}
\begin{adjustwidth}{-\extralength}{0cm}
\centering
\begin{tabular}{lccccc}
\toprule
\textbf{Channel} & \textbf{Mesh (mm)} & $L$ \textbf{(DC)} & $L$ \textbf{(\SI{1}{\mega\hertz})} & \textbf{Loop len.} & \textbf{Eff. width} \\
\midrule
Fuel cell & 0.3 & \SI{1.538}{\nano\henry} & \SI{1.456}{\nano\henry} & \SI{4.74}{\milli\meter} & \SI{12.47}{\milli\meter} \\
Battery   & 0.2 & \SI{3.480}{\nano\henry} & \SI{3.149}{\nano\henry} & \SI{9.69}{\milli\meter} & \SI{5.91}{\milli\meter} \\
\bottomrule
\end{tabular}
\end{adjustwidth}
\end{table}

% NOTE: the fix + validation. The single-variable framing is the strongest
% claim in this whole section; state the controls explicitly.

The indicated remedy is a board-level one---firmware cannot reduce a loop
inductance. Ceramic capacitors of \SI{10}{\micro\farad} and \SI{0.1}{\micro\farad}
were bodged directly at the battery converter's output pins, collapsing the
\SI{240}{mil} loop to approximately the fuel-cell channel's geometry. The
validation was deliberately single-variable: an otherwise exact repetition of the
Death~4 configuration, differing only in the added capacitors. Where every
unmodified attempt had destroyed a converter, the modified channel survived
\emph{four consecutive} bring-ups, regulating the bus normally
(Figures~\ref{fig:scope-vout}--\ref{fig:scope-sw}).

% NOTE: this caveat is load-bearing for honesty. Do not cut it in editing.
Two limitations qualify this result. First, all validation captures were taken at
$1\times$ probe bandwidth ($\approx$\SI{10}{\mega\hertz}, \SI{50}{MSa/s}); the
hypothesized \SIrange{100}{200}{\mega\hertz} hot-loop ring is invisible at that
bandwidth. \emph{Survival is the evidence, not the waveform}, and a high-bandwidth
switch-node margin measurement remains owed. Second, the extraction models one
user-selected loop at a time and is not a full-board three-dimensional solve; it
supports a comparative claim between two nominally identical channels far better
than it supports an absolute ring-amplitude prediction.

% ------------------------------------------------------------------
\subsection{Soft-Start and Protection Timing in the Ideal-Diode Path}
\label{sec:bringup:softstart}

% NOTE: this is the subsection to sell to a *control* audience. Frame it as two
% independently designed loops (a fixed-timing open-loop charge ramp and a
% closed-loop voltage regulator) coupled through a node neither one models.

The hot-loop fix held for the battery channel, but a fifth converter---this time
on the fuel-cell side, with good layout---was destroyed while closing the motor-node
switch onto an attached motor controller. This failure has a different character
and is, in our view, the more generalizable one, because it arises from the
\emph{interaction of two independently designed control loops} rather than from a
layout defect.

The bus-to-motor path is gated by an RT1987 ideal-diode controller whose inrush
behavior is set by an open-loop soft-start ramp: a fixed external capacitance
($C_\mathrm{SS}=\SI{5.6}{\nano\farad}$) sets a turn-on time of order
\SI{1.1}{\milli\second}, during which a start-up short-circuit protection (SCP)
foldback current limit is armed with a \SI{250}{\micro\second} continuous-clamp
timer and a \SI{64}{\milli\second} retry interval. The ramp time is a designed
constant; it does not adapt to the load it is charging. Downstream sits the motor
node---\SI{470}{\micro\farad} of bulk capacitance plus the motor controller's own
(unmeasured) input capacitance, a total measured near \SI{590}{\micro\farad} on the
bare board. The charge demand of that node over the fixed ramp,
$I \approx 0.8\,C\,V/t_\mathrm{ON}$, lands in the same amperes-class range as the
applicable foldback limit. The ideal diode's protection therefore engages or nearly
engages on connection, and each clamp-and-release cycle presents the upstream boost
converters with a large, abrupt load step followed by an equally abrupt load dump.

% NOTE: the control-theoretic reading. This is the sentence the section exists for.
Neither loop is misdesigned in isolation. The ideal diode's soft start is a
correct open-loop inrush limiter for a bounded load capacitance; the boost's
voltage loop is a correctly compensated regulator with a \SIrange{4}{19}{\kilo\hertz}
crossover. The hazard lives in their coupling: the load-dump transient at the end
of a clamp interval occurs on a timescale comparable to the boost's loop latency,
so the boost error amplifier is left wound up delivering current into a suddenly
unloaded output. The resulting overshoot is then \emph{peak-held} by the ideal
diode itself, which cannot sink current---the switch behaves as a peak detector on
the bus, and the rail parks above nominal until leakage discharges it.
Measured parks of \SIrange{0.65}{0.72}{\volt} (centre-to-centre) above the
regulation point were observed, sufficient to trip the firmware bus overvoltage
fault on roughly \SI{80}{\percent} of bring-up attempts once the bus nominal was
retuned downward. At the \SI{15}{\ampere}-class currents of a full-bus motor-node
connection, the same event scaled past the converter's absolute-maximum rating and
destroyed it---but only from a stiff source, since a battery source sags into
undervoltage lockout before lethal current flows. This source dependence
retroactively explains the otherwise puzzling pattern across all five failures.

% TODO(verify): the log flags the *specific* clamp/cut identification as
% UNCONFIRMED and partially superseded -- the deepest observed event is now
% believed to be a soft-start that COMPLETED (1.77 ms conduction, current
% tapering as delta-V goes to zero) rather than a timer cut, and a smaller
% earlier event violates charge conservation by 14-300x and may be a rectified
% ring artifact of the unipolar current sensor. Write this subsection so that
% the *architectural* claim (fixed-ramp open-loop soft start vs. closed-loop
% regulator, coupled through an unmodeled capacitive node) survives regardless
% of which micro-mechanism wins at the bench. Do NOT publish the micro-mechanism
% as settled.

Two remedies follow directly, and are complementary. The firmware remedy is
\emph{sequencing}: never close the motor-node switch at full bus voltage. The
bring-up now raises the motor-node switch together with the bus switches
\emph{before} the boost converters ramp, so the converters' own soft start charges
the entire downstream stack from a low initial voltage in one coordinated ramp; the
node is then kept energized through idle and run states, with the motor held
stopped by a zero current command rather than by removing power. A guard predicate
faults rather than permitting a hot connection if the bus and motor-node
measurements indicate the node is discharged. This trades away hardware motor
isolation in idle---a deliberate and documented posture change. The hardware remedy
is to size the ideal diode's soft-start capacitance to the load it actually
charges: increasing $C_\mathrm{SS}$ from \SI{5.6}{\nano\farad} to
\SI{100}{\nano\farad} extends the ramp to tens of milliseconds and reduces the
connect demand to a few hundred milliamperes, restoring roughly an order of
magnitude of margin. A subsequent datasheet reconciliation also revealed that the
firmware's inter-phase settle delay had been sized from the soft-start time while
omitting the controller's \SI{8}{\milli\second} turn-on \emph{delay}, so the
intended coordinated ramp was not in fact executing; the current recommendation is
to replace the fixed delay with an ADC-confirmed pre-charge gate and timeout rather
than any tuned constant.

% ------------------------------------------------------------------
\subsection{Relationship to the Droop Architecture}
\label{sec:bringup:droop-relevance}

% NOTE: reviewers WILL ask whether droop caused any of this. Answer per-issue,
% in a table, and concede the one place where droop has genuine collateral
% (R_C is load-bearing for the share plant).

It is important to state precisely which of these hazards the droop power-sharing
architecture created, and which it merely inherited. Table~\ref{tab:droop-relevance}
gives the attribution.

\begin{table}[H]
\caption{Attribution of each bring-up hazard to architecture-specific versus generic causes.\label{tab:droop-relevance}}
\begin{adjustwidth}{-\extralength}{0cm}
\centering
\begin{tabular}{p{4.2cm}lp{7.5cm}}
\toprule
\textbf{Hazard} & \textbf{Attribution} & \textbf{Basis} \\
\midrule
Output hot-loop overshoot (Deaths 1--4) & Generic & A single-channel boost driving any capacitive bus exhibits it. The droop injection network was inactive (DACs unwritten) in every failure; the asymmetry is fabrication layout, present before any control code ran. \\
Motor-node hot-plug load dump (Death 5) & Generic multi-converter & Arises from an ideal-diode soft-start versus a downstream capacitive node; identical in a single-source system with the same power-path switching. \\
Ideal-diode peak-hold overvoltage parking & Generic & A consequence of non-sinking sources behind unidirectional switches; independent of how current is shared between them. \\
Source-collapse restart cycling & Generic (bench) & An artifact of deriving logic power from a converter input rail on a soft supply. \\
Droop-injection gain mapping error & \textbf{Droop-specific} & A firmware mapping omitted the feedback-injection attenuation, saturating both DACs and presenting the share loop a zero-gain plant. Non-destructive, but it invalidated all share-loop bench data taken before it was found (Section~\ref{sec:droop:bug}). \\
Compensation-resistor lever & \textbf{Coupled} & The obvious mitigation for load-dump overshoot---raising the boost crossover---is constrained because the compensation resistor sets the response time constant of the share-loop plant model, so a unilateral change would invalidate the synthesized controller. \\
\bottomrule
\end{tabular}
\end{adjustwidth}
\end{table}

The summary judgement is that the destructive failures were \emph{multi-converter
power-path bring-up hazards}, not droop-control hazards: they occurred with the
sharing loop open, with the modulating DACs at their reset state, and on a
single-channel basis. The droop architecture's genuine contribution to the campaign
was a non-destructive but consequential one---a gain-mapping error that silently
opened the share loop---and a constraint on the remedy space, in that the
compensation network is shared between the local voltage loop and the outer sharing
plant and cannot be retuned for transient margin without re-synthesizing the
controller. That coupling is, we would argue, the transferable control-level lesson:
in a droop-shared converter pair, the per-converter compensation network is no longer
a purely local design variable.

% ------------------------------------------------------------------
\subsection{Figures}
\label{sec:bringup:figures}

% NOTE: all placeholders. Scope JPGs exist under references/scope_captures/ and
% need re-export at publication resolution with cursors legible; the two drawn
% figures do not exist yet.

\begin{figure}[H]
\centering
% \includegraphics[width=0.8\textwidth]{Figures/scope-vout-validation.pdf}
\fbox{\parbox[c][4cm][c]{0.8\textwidth}{\centering PLACEHOLDER --- source:
\texttt{references/scope\_captures/1-VOUT.jpg}}}
\caption{Battery-converter output during a validated bring-up with hot-loop
capacitors fitted: body-diode pre-charge, one aborted soft-start attempt with
automatic retry, then clean regulation. The retry is a benign source-sag
undervoltage event, and is the same class of event that destroyed converters
before the fix.\label{fig:scope-vout}}
\end{figure}

\begin{figure}[H]
\centering
% \includegraphics[width=0.8\textwidth]{Figures/scope-triple.pdf}
\fbox{\parbox[c][4cm][c]{0.8\textwidth}{\centering PLACEHOLDER --- source:
\texttt{references/scope\_captures/5-Vout yellow-Vmot blue-Ibat purple.jpg}}}
\caption{Three-channel capture through a bring-up: converter output, bus voltage,
and current-sense output. The wide current hump is the downstream node charging
through a soft start that completes; the subsequent output park is the peak-held
overshoot.\label{fig:scope-triple}}
% TODO: cursor values in the raw capture are readable only at full resolution;
% re-export and annotate. Confirm the axis scaling caption (5 V/div, 5 A/div,
% 2 ms/div) against the instrument, not the log.
\end{figure}

\begin{figure}[H]
\centering
% \includegraphics[width=0.7\textwidth]{Figures/scope-sw-zoom.pdf}
\fbox{\parbox[c][3.5cm][c]{0.7\textwidth}{\centering PLACEHOLDER --- source:
\texttt{references/scope\_captures/3-SW.jpg}, \texttt{4-SW zoomed in.jpg}}}
\caption{Switch-node envelope during soft start. The initial discrete pulses are
normal pulse-frequency-modulation skipping, not instability. Note the bandwidth
limitation discussed in Section~\ref{sec:bringup:inductance}.\label{fig:scope-sw}}
\end{figure}

\begin{figure}[H]
\centering
% \includegraphics[width=0.8\textwidth]{Figures/inductance-sweep.pdf}
\fbox{\parbox[c][5cm][c]{0.8\textwidth}{\centering PLACEHOLDER --- generate from
\texttt{tools/inductance} frequency-sweep store}}
\caption{Extracted loop inductance versus frequency for both converter output
loops, with the mesh-convergence axis inset and the closed-form microstrip bound
shown for reference.\label{fig:Lsweep}}
% TODO: two-panel figure. (a) L(f) DC-1 MHz, both channels. (b) L vs mesh pitch
% at the designated frequency, showing convergence. Overlay the analytic bound
% as a dashed line -- it sits well below the solved value (0.64 vs 1.54 nH FC;
% 2.37 vs 3.48 nH BT) and that gap needs one sentence of explanation in the
% caption, since a reviewer will otherwise read it as a validation failure.
\end{figure}

\begin{figure}[H]
\centering
% \includegraphics[width=\textwidth]{Figures/fault-tree.pdf}
\fbox{\parbox[c][6cm][c]{\textwidth}{\centering PLACEHOLDER --- to be drawn}}
\caption{Campaign timeline and fault tree. Upper band: chronological failure and
intervention events. Lower band: hypothesis tree, with each branch annotated by the
experiment that refuted or confirmed it.\label{fig:faulttree}}
% TODO: draw in TikZ so it stays vector and editable. Suggested structure --
% root "converter destroyed on bus connect"; first-level branches: part defect,
% supply/source, bus path fault, firmware sequencing, channel layout; each leaf
% carrying its refuting experiment. Mark the surviving branch and the point at
% which Death 5 forced a second, independent root cause to be opened.
\end{figure}


% ============================================================================
% PLANNING NOTES
% ============================================================================
%
% PURPOSE OF THIS SECTION IN THE PAPER
% - Serves the "systems modeling" half of the thesis claim. The control chapters
%   show a model-based synthesis; this chapter shows model-based *diagnosis*.
%   Both leaned on building a quantitative model when iteration had stalled.
% - Secondary purpose: pre-empt the reviewer question "did your droop scheme
%   destroy your hardware?" The attribution table answers it explicitly and
%   concedes the one real coupling (compensation network shared between local
%   and sharing loops).
%
% OPEN ITEMS BEFORE THIS SECTION CAN BE SUBMITTED
% 1. INDUCTANCE RATIO DISCREPANCY. Project docs quote ~2.7x; the current results
%    store gives 2.16x @ 1 MHz. The discrepancy is a loop-closure definition
%    change (pre-2026-07-01 numbers are stale per project memory). Re-run the
%    sweep with the current closure on both configs, pick one number, and purge
%    the other from the paper, CLAUDE.md, and the debug log simultaneously.
% 2. ANALYTIC-BOUND GAP. The closed-form bound is ~2.4x below the solved value on
%    both channels. Consistent bias, so it does not threaten the comparative
%    claim, but a reviewer will ask. Either explain (bound neglects return-path
%    spreading / via inductance) or drop the bound from the figure entirely.
% 3. HIGH-BANDWIDTH RING MEASUREMENT IS STILL OWED. The entire hot-loop mechanism
%    rests on an inferred 100-200 MHz ring that has never been observed -- every
%    capture to date is ~10 MHz bandwidth. Until a 10x probe with a ground spring
%    captures the switch node at a controlled load step, the mechanism is
%    "consistent with all evidence", not "measured". Language throughout this
%    section must reflect that. Currently it does; keep it that way.
% 4. DEATH-5 MICRO-MECHANISM IS IN FLUX. The debug log has superseded its own
%    SCP-cut interpretation once already (a capture showed the soft start
%    completing rather than being cut). The architectural framing used here is
%    robust to that, but check the log's current state immediately before
%    submission and re-verify every number quoted in that subsection.
% 5. VESC INPUT CAPACITANCE IS UNMEASURED anywhere in the project. It appears in
%    the C_SS sizing argument as an unbounded term. Measure it, or drop the
%    quantitative margin claim to a qualitative one.
% 6. Deaths 1-4 count as "four battery-side boosts" in the summary doc but the
%    combined narrative with Death 5 makes five total. Make sure the abstract and
%    intro use the same count.
%
% STRUCTURAL DECISIONS TO REVISIT
% - The extraction-tool subsection may deserve promotion to its own methods
%   section or even a standalone short paper -- it is a reusable Gerber-to-PEEC
%   pipeline and is the most novel artifact here. Decide after the control
%   chapters are drafted and total length is known.
% - Consider moving the deaths table to an appendix and keeping only the
%   narrative, if Energies page pressure bites. The fault tree figure can carry
%   the chronology.
% - Tone check: the failure narrative is deliberately unflattering. That is the
%   point (superseded-hypothesis reporting is the contribution), but the
%   introduction to the chapter must frame it as method, not confession, or a
%   reviewer will read it as inexperience.
%
% CITATIONS TO ADD
% - TPS61288 datasheet (abs-max, OVP thresholds, compensation guidance §9.2.2.5)
% - RT1987 datasheet (tD_ON, soft-start timing, SCP foldback)
% - INA253 datasheet (unipolar sensing -- relevant to the half-rectified ring
%   caveat if that caveat survives into the final draft)
% - FastHenry / FastFieldSolvers, and a PEEC reference (Ruehli) for the method
% - A hot-loop layout reference (TI app notes on buck/boost layout) to
%   support the "Cout distance is the metric, not trace length" claim
% ============================================================================
